\textbf{Phần 1. Trắc nghiệm lựa chọn}
\setcounterpageref{ex}{0}

\begin{ex} %[5H5N29]
Một hình tròn có bán kính $50\text{ cm}$ và một hình tròn khác có bán kính $200\text{ cm}$. Chu vi hình tròn lớn gấp chu vi hình tròn bé số lần là
\choice
{\True $4$ lần}
{$2$ lần}
{$3$ lần}
{$5$ lần}
\loigiai{
Tỉ số bán kính hai hình tròn là $200 : 50 = 4$. Do đó chu vi hình tròn lớn gấp $4$ lần chu vi hình tròn bé.\\
\textbf{Đáp án của câu hỏi là A}
}
\end{ex}

\begin{ex} %[5H5H29]
Một cái ao có dạng nửa hình tròn đường kính $12\text{ m}$. Chu vi của cái ao đó (gồm đường cong nửa hình tròn và đường kính bờ ao) là
\choice
{$18{,}84\text{ m}$}
{$37{,}68\text{ m}$}
{\True $30{,}84\text{ m}$}
{$49{,}68\text{ m}$}
\loigiai{
Độ dài đường cong nửa hình tròn là: $3{,}14 \times 12 : 2 = 18{,}84\text{ (m)}$.\\
Chu vi của cái ao là: $18{,}84 + 12 = 30{,}84\text{ (m)}$.\\
\textbf{Đáp án của câu hỏi là C}
}
\end{ex}

\begin{ex} %[5H5H29]
Cho hình chữ nhật $ABCD$ có chiều dài $b\text{ cm}$ và chiều rộng $a\text{ cm}$. Lấy điểm $E$ bất kì trên cạnh $AB$. Diện tích hình tam giác $ECD$ là
\choice
{\True $\dfrac{a \times b}{2}\text{ (cm}^2\text{)}$}
{$a \times b\text{ (cm}^2\text{)}$}
{$\dfrac{a \times b}{4}\text{ (cm}^2\text{)}$}
{Thay đổi tùy thuộc vào vị trí điểm $E$}
\loigiai{
Tam giác $ECD$ có đáy là $CD = b\text{ cm}$ và chiều cao hạ từ $E$ xuống $CD$ bằng chiều rộng $a\text{ cm}$. Do đó diện tích tam giác $ECD$ luôn bằng $\dfrac{a \times b}{2}\text{ (cm}^2\text{)}$.\\
\textbf{Đáp án của câu hỏi là A}
}
\end{ex}

\begin{ex} %[5H5H29]
Một hình vuông có cạnh $8\text{ cm}$ chứa một hình tròn có đường kính $8\text{ cm}$ bên trong. Diện tích phần hình vuông nằm ngoài hình tròn là
\choice
{$50{,}24\text{ cm}^2$}
{$114{,}24\text{ cm}^2$}
{$136{,}96\text{ cm}^2$}
{\True $13{,}76\text{ cm}^2$}
\loigiai{
Diện tích hình vuông là: $8 \times 8 = 64\text{ (cm}^2\text{)}$.\\
Bán kính hình tròn là: $8 : 2 = 4\text{ (cm)}$.\\
Diện tích hình tròn là: $3{,}14 \times 4 \times 4 = 50{,}24\text{ (cm}^2\text{)}$.\\
Diện tích phần màu xanh là: $64 - 50{,}24 = 13{,}76\text{ (cm}^2\text{)}$.\\
\textbf{Đáp án của câu hỏi là D}
}
\end{ex}

\begin{ex} %[5H5H29]
Bạn Việt nặn ba hình từ đất sét: hình tam giác vuông có hai cạnh góc vuông dài $7\text{ cm}$ và $7\text{ cm}$; hình thang có hai đáy $3\text{ cm}, 6\text{ cm}$ và chiều cao $4\text{ cm}$; hình tròn bán kính $5\text{ cm}$. Hình có diện tích bé nhất là
\choice
{Hình tam giác}
{Hình tròn}
{\True Hình thang}
{Ba hình có diện tích bằng nhau}
\loigiai{
Diện tích hình tam giác: $\dfrac{7 \times 7}{2} = 24{,}5\text{ cm}^2$.\\
Diện tích hình thang: $\dfrac{(3 + 6) \times 4}{2} = 18\text{ cm}^2$.\\
Diện tích hình tròn: $3{,}14 \times 5 \times 5 = 78{,}5\text{ cm}^2$.\\
Vậy hình thang có diện tích bé nhất ($18\text{ cm}^2$).\\
\textbf{Đáp án của câu hỏi là C}
}
\end{ex}

\begin{ex} %[5H5H29]
Trong ba hình ở câu trên, hình có diện tích lớn nhất là
\choice
{Hình tam giác}
{\True Hình tròn}
{Hình thang}
{Không có hình lớn nhất}
\loigiai{
Hình tròn có diện tích lớn nhất ($78{,}5\text{ cm}^2$).\\
\textbf{Đáp án của câu hỏi là B}
}
\end{ex}

\begin{ex} %[5H5V29]
Một mảnh đất gồm hình thang $ABCD$ (hai đáy $30\text{ m}$ và $55\text{ m}$, chiều cao $25\text{ m}$) và hình tam giác $AED$ (đáy $55\text{ m}$, chiều cao $27\text{ m}$). Diện tích toàn bộ mảnh đất đó là
\choice
{$1\,062{,}5\text{ m}^2$}
{$742{,}5\text{ m}^2$}
{$1\,500\text{ m}^2$}
{\True $1\,805\text{ m}^2$}
\loigiai{
Diện tích hình thang $ABCD$: $\dfrac{(55 + 30) \times 25}{2} = 1\,062{,}5\text{ (m}^2\text{)}$.\\
Diện tích tam giác $AED$: $\dfrac{55 \times 27}{2} = 742{,}5\text{ (m}^2\text{)}$.\\
Diện tích toàn bộ mảnh đất: $1\,062{,}5 + 742{,}5 = 1\,805\text{ (m}^2\text{)}$.\\
\textbf{Đáp án của câu hỏi là D}
}
\end{ex}

\begin{ex} %[5H5V29]
Cho hình vuông $ABCD$ cạnh $6{,}5\text{ cm}$. Trên cạnh $CD$ lấy các điểm sao cho $DE = EG = GH = HK = KC = 1{,}3\text{ cm}$. Diện tích hình thang vuông $ABCK$ là
\choice
{$50{,}7\text{ cm}^2$}
{\True $25{,}35\text{ cm}^2$}
{$12{,}675\text{ cm}^2$}
{$42{,}25\text{ cm}^2$}
\loigiai{
Đáy $AB = 6{,}5\text{ cm}$, đáy $KC = 1{,}3\text{ cm}$, chiều cao $BC = 6{,}5\text{ cm}$.\\
Diện tích hình thang $ABCK$ là: $\dfrac{(6{,}5 + 1{,}3) \times 6{,}5}{2} = \dfrac{7{,}8 \times 6{,}5}{2} = 25{,}35\text{ (cm}^2\text{)}$.\\
\textbf{Đáp án của câu hỏi là B}
}
\end{ex}

\begin{ex} %[5H5V29]
Cũng trong hình vuông $ABCD$ ở câu trên, diện tích tam giác $AKD$ gấp diện tích tam giác $ADE$ số lần là
\choice
{$2$ lần}
{$3$ lần}
{\True $4$ lần}
{$5$ lần}
\loigiai{
Đáy $KD = 4 \times 1{,}3 = 5{,}2\text{ cm}$ gấp $4$ lần đáy $DE = 1{,}3\text{ cm}$. Hai tam giác có cùng chiều cao $AD$ nên diện tích tam giác $AKD$ gấp $4$ lần diện tích tam giác $ADE$.\\
\textbf{Đáp án của câu hỏi là C}
}
\end{ex}

\begin{ex} %[5H5V29]
Một chiếc khăn trải bàn hình chữ nhật có chiều dài $2\text{ m}$ và chiều rộng $1{,}5\text{ m}$. Ở giữa khăn người ta thêu một hình thoi có các đường chéo dài $1\text{ m}$ và $0{,}8\text{ m}$. Diện tích phần vải không thêu là
\choice
{$3\text{ m}^2$}
{\True $2{,}6\text{ m}^2$}
{$0{,}4\text{ m}^2$}
{$2{,}2\text{ m}^2$}
\loigiai{
Diện tích khăn trải bàn: $2 \times 1{,}5 = 3\text{ (m}^2\text{)}$.\\
Diện tích hình thoi: $\dfrac{1 \times 0{,}8}{2} = 0{,}4\text{ (m}^2\text{)}$.\\
Diện tích phần không thêu: $3 - 0{,}4 = 2{,}6\text{ (m}^2\text{)}$.\\
\textbf{Đáp án của câu hỏi là B}
}
\end{ex}

\begin{ex} %[5H5V29]
Một bánh xe đạp có bán kính $0{,}35\text{ m}$. Khi người đi xe đạp đi được quãng đường $219{,}8\text{ m}$ thì bánh xe đó đã lăn được số vòng là
\choice
{$50$ vòng}
{$150$ vòng}
{$200$ vòng}
{\True $100$ vòng}
\loigiai{
Chu vi bánh xe là: $3{,}14 \times 2 \times 0{,}35 = 2{,}198\text{ (m)}$.\\
Số vòng lăn của bánh xe là: $219{,}8 : 2{,}198 = 100$ (vòng).\\
\textbf{Đáp án của câu hỏi là D}
}
\end{ex}

\begin{ex} %[5H5V29]
Nếu bán kính của một hình tròn tăng thêm $20\%$ thì diện tích của hình tròn đó tăng thêm bao nhiêu phần trăm?
\choice
{\True $44\%$}
{$20\%$}
{$40\%$}
{$48\%$}
\loigiai{
Bán kính mới bằng $120\% = 1{,}2$ lần bán kính cũ.\\
Diện tích mới bằng: $1{,}2 \times 1{,}2 = 1{,}44 = 144\%$ diện tích cũ.\\
Diện tích tăng thêm: $144\% - 100\% = 44\%$.\\
\textbf{Đáp án của câu hỏi là A}
}
\end{ex}

\textbf{Phần 2. Trắc nghiệm đúng sai}
\setcounterpageref{ex}{0}

\begin{ex} %[5H5H29]
Cho hình tròn đỏ có đường kính gấp $2$ lần đường kính của hình tròn xanh. Xét tính đúng sai của các khẳng định sau:
\choiceTF
{\True Bán kính hình tròn đỏ gấp $2$ lần bán kính hình tròn xanh.}
{\True Chu vi hình tròn đỏ gấp $2$ lần chu vi hình tròn xanh.}
{Diện tích hình tròn đỏ gấp $2$ lần diện tích hình tròn xanh.}
{\True Diện tích hình tròn đỏ gấp $4$ lần diện tích hình tròn xanh.}
\loigiai{
\begin{itemize}
    \item $d_{đỏ} = 2 \times d_{xanh} \Rightarrow r_{đỏ} = 2 \times r_{xanh}$.
    \item Chu vi: $C_{đỏ} = 3{,}14 \times d_{đỏ} = 2 \times C_{xanh}$.
    \item Diện tích: $S_{đỏ} = 3{,}14 \times (2 \times r_{xanh}) \times (2 \times r_{xanh}) = 4 \times S_{xanh} \neq 2 \times S_{xanh}$.
    \item Diện tích hình tròn đỏ gấp 4 lần diện tích hình tròn xanh.
\end{itemize}
\begin{itemize}
    \item a) \textbf{Đúng}.
    \item b) \textbf{Đúng}.
    \item c) \textbf{Sai}.
    \item d) \textbf{Đúng}.
\end{itemize}
}
\end{ex}

\begin{ex} %[5H5V29]
Một mảnh đất có dạng như hình vẽ được chia thành hình thang $ABCD$ (đáy $30\text{ m}$ và $55\text{ m}$, chiều cao $25\text{ m}$) và hình tam giác $AED$ (đáy $55\text{ m}$, chiều cao $27\text{ m}$). Xét tính đúng sai của các khẳng định sau:
\choiceTF
{\True Diện tích phần hình thang $ABCD$ là $1\,062{,}5\text{ m}^2$.}
{\True Diện tích phần hình tam giác $AED$ là $742{,}5\text{ m}^2$.}
{Diện tích phần tam giác $AED$ lớn hơn diện tích phần hình thang $ABCD$.}
{\True Diện tích toàn bộ mảnh đất là $1\,805\text{ m}^2$.}
\loigiai{
\begin{itemize}
    \item Diện tích hình thang: $\dfrac{(55 + 30) \times 25}{2} = 1\,062{,}5\text{ m}^2$.
    \item Diện tích hình tam giác: $\dfrac{55 \times 27}{2} = 742{,}5\text{ m}^2$.
    \item $742{,}5\text{ m}^2 < 1\,062{,}5\text{ m}^2$ nên diện tích tam giác bé hơn diện tích hình thang.
    \item Tổng diện tích: $1\,062{,}5 + 742{,}5 = 1\,805\text{ m}^2$.
\end{itemize}
\begin{itemize}
    \item a) \textbf{Đúng}.
    \item b) \textbf{Đúng}.
    \item c) \textbf{Sai}.
    \item d) \textbf{Đúng}.
\end{itemize}
}
\end{ex}

\textbf{Phần 3. Trả lời ngắn}
\setcounterpageref{ex}{0}

\begin{ex} %[5H5H29]
Cái ao dạng nửa hình tròn đường kính $12\text{ m}$. Tính chu vi của cái ao đó theo mét. (Nhập kết quả dưới dạng số thập phân).
\shortans{30,84}
\loigiai{
\textbf{Đáp số:} $30{,}84$\\
$3{,}14 \times 12 : 2 + 12 = 18{,}84 + 12 = 30{,}84\text{ (m)}$.
}
\end{ex}

\begin{ex} %[5H5H29]
Hình vuông cạnh $8\text{ cm}$ chứa hình tròn đường kính $8\text{ cm}$. Tính diện tích phần hình vuông nằm ngoài hình tròn theo xăng-ti-mét vuông. (Nhập kết quả dưới dạng số thập phân).
\shortans{13,76}
\loigiai{
\textbf{Đáp số:} $13{,}76$\\
$8 \times 8 - 3{,}14 \times 4 \times 4 = 64 - 50{,}24 = 13{,}76\text{ (cm}^2\text{)}$.
}
\end{ex}

\begin{ex} %[5H5V29]
Hình thang có hai đáy $30\text{ m}, 55\text{ m}$, chiều cao $25\text{ m}$ và tam giác có đáy $55\text{ m}$, chiều cao $27\text{ m}$. Tính tổng diện tích của hai hình theo mét vuông.
\shortans{1805}
\loigiai{
\textbf{Đáp số:} $1805$\\
$\dfrac{(55 + 30) \times 25}{2} + \dfrac{55 \times 27}{2} = 1\,062{,}5 + 742{,}5 = 1\,805\text{ (m}^2\text{)}$.
}
\end{ex}

\begin{ex} %[5H5V29]
Hình vuông cạnh $6{,}5\text{ cm}$. Tính diện tích hình thang $ABCK$ có hai đáy $6{,}5\text{ cm}$ và $1{,}3\text{ cm}$, chiều cao $6{,}5\text{ cm}$ theo xăng-ti-mét vuông. (Nhập kết quả dưới dạng số thập phân).
\shortans{25,35}
\loigiai{
\textbf{Đáp số:} $25{,}35$\\
$\dfrac{(6{,}5 + 1{,}3) \times 6{,}5}{2} = 25{,}35\text{ (cm}^2\text{)}$.
}
\end{ex}

\begin{ex} %[5H5V29]
Một bánh xe đạp có bán kính $0{,}35\text{ m}$. Hỏi khi bánh xe lăn được $100$ vòng thì xe đi được quãng đường bao nhiêu mét? (Nhập kết quả dưới dạng số thập phân).
\shortans{219,8}
\loigiai{
\textbf{Đáp số:} $219{,}8$\\
$3{,}14 \times 2 \times 0{,}35 \times 100 = 219{,}8\text{ (m)}$.
}
\end{ex}

\begin{ex} %[5H5V29]
Một mảnh khăn hình chữ nhật $2\text{ m} \times 1{,}5\text{ m}$ thêu hình thoi ở giữa có hai đường chéo dài $1\text{ m}$ và $0{,}8\text{ m}$. Tính diện tích phần khăn không thêu theo mét vuông. (Nhập kết quả dưới dạng số thập phân).
\shortans{2,6}
\loigiai{
\textbf{Đáp số:} $2{,}6$\\
$2 \times 1{,}5 - \dfrac{1 \times 0{,}8}{2} = 3 - 0{,}4 = 2{,}6\text{ (m}^2\text{)}$.
}
\end{ex}
